\documentclass[11pt]{report}
% Same style as 1-foundations/theory/preamble.tex (the course's own documents).

\usepackage[a4paper,margin=2.3cm,headheight=14pt]{geometry}
\usepackage[T1]{fontenc}
\usepackage[utf8]{inputenc}
\usepackage{lmodern}
\usepackage{microtype}
\usepackage{booktabs}
\usepackage{longtable}
\usepackage{array}
\usepackage{xcolor}
\usepackage{enumitem}
\usepackage{amssymb}
\usepackage{amsmath}
\usepackage{fancyhdr}
\usepackage{listings}
\usepackage{tikz}
\usetikzlibrary{arrows.meta,positioning,calc,fit,shapes.geometric,decorations.pathreplacing}
\usepackage[hidelinks]{hyperref}

\definecolor{warnbg}{RGB}{255,247,230}
\definecolor{warnrule}{RGB}{204,136,0}
\definecolor{newbg}{RGB}{235,245,255}
\definecolor{newrule}{RGB}{40,110,190}
\definecolor{keybg}{RGB}{237,247,237}
\definecolor{keyrule}{RGB}{50,120,60}
\definecolor{headgray}{RGB}{90,90,90}
\definecolor{codebg}{RGB}{248,248,248}

\usepackage[most]{tcolorbox}

\newtcolorbox{warnbox}[1][]{colback=warnbg,colframe=warnrule,boxrule=0.8pt,
  left=6pt,right=6pt,top=5pt,bottom=5pt,arc=1pt,#1}
\newtcolorbox{infobox}[1][]{colback=newbg,colframe=newrule,boxrule=0.8pt,
  left=6pt,right=6pt,top=5pt,bottom=5pt,arc=1pt,#1}
% "In SWEB": where the idea shows up in the course's kernel
\newtcolorbox{swebbox}[1][]{colback=newbg,colframe=newrule,boxrule=0.8pt,
  left=6pt,right=6pt,top=5pt,bottom=5pt,arc=1pt,
  title={In SWEB},fonttitle=\bfseries\small,coltitle=white,colbacktitle=newrule,#1}
% end-of-chapter summary
\newtcolorbox{keybox}[1][]{colback=keybg,colframe=keyrule,boxrule=0.8pt,
  left=6pt,right=6pt,top=5pt,bottom=5pt,arc=1pt,
  title={Key points},fonttitle=\bfseries\small,coltitle=white,colbacktitle=keyrule,#1}

\newcommand{\tbd}{\textcolor{warnrule}{\textbf{[confirm]}}}
\newcommand{\code}[1]{\texttt{#1}}
\newcommand{\file}[1]{\texttt{\small #1}}
\newcommand{\module}[1]{\textsf{module~#1}}

\setlist[itemize]{topsep=2pt,itemsep=1.5pt,parsep=0pt,leftmargin=1.3em}
\setlist[enumerate]{topsep=2pt,itemsep=1.5pt,parsep=0pt,leftmargin=1.6em}
\setlist[description]{topsep=2pt,itemsep=1.5pt,parsep=0pt,leftmargin=1.3em}

\newcolumntype{L}[1]{>{\raggedright\arraybackslash}p{#1}}

\lstset{
  language=C,
  basicstyle=\ttfamily\small,
  keywordstyle=\color{newrule}\bfseries,
  commentstyle=\color{headgray}\itshape,
  stringstyle=\color{warnrule!80!black},
  backgroundcolor=\color{codebg},
  frame=single, rulecolor=\color{black!15}, framesep=3pt,
  xleftmargin=4pt, xrightmargin=4pt, aboveskip=6pt, belowskip=4pt,
  columns=fullflexible, keepspaces=true, showstringspaces=false,
  upquote=true, tabsize=2,
  morekeywords={size_t,ssize_t,pid_t,pthread_t,pthread_mutex_t,pthread_cond_t,
                sem_t,uint8_t,uint32_t,uint64_t,bool,_Atomic,atomic_int,
                atomic_flag,__thread,_Noreturn,__asm__,volatile,pointer,uint64}
}
\lstdefinestyle{shell}{language=bash,morekeywords={},keywordstyle=\color{black}}
\lstdefinestyle{cpp}{language=C++}

\pagestyle{fancy}
\fancyhf{}
\renewcommand{\headrulewidth}{0.4pt}
\fancyhead[L]{\small\textcolor{headgray}{SWEB interview guide}}
\fancyhead[R]{\small\textcolor{headgray}{\leftmark}}
\fancyfoot[C]{\small\thepage}
\renewcommand{\chaptermark}[1]{\markboth{\thechapter\ \ #1}{}}
\fancypagestyle{plain}{\fancyhf{}\renewcommand{\headrulewidth}{0pt}\fancyfoot[C]{\small\thepage}}

\renewcommand{\arraystretch}{1.25}
\setlength{\parskip}{3pt}
\setlength{\parindent}{0pt}
\setlength{\emergencystretch}{3em}   % fewer overfull lines with long \code words

% guide additions
\newtcolorbox{lockbox}[1][]{colback=warnbg,colframe=warnrule,boxrule=0.8pt,
  left=6pt,right=6pt,top=5pt,bottom=5pt,arc=1pt,
  title={Locking},fonttitle=\bfseries\small,coltitle=white,colbacktitle=warnrule,#1}
\newtcolorbox{trapbox}[1][]{colback=warnbg,colframe=warnrule,boxrule=0.8pt,
  left=6pt,right=6pt,top=5pt,bottom=5pt,arc=1pt,
  title={Trap},fonttitle=\bfseries\small,coltitle=white,colbacktitle=warnrule,#1}
\lstset{morekeywords={ScopeLock,Mutex,Condition,SpinLock,ArchThreadRegisters,ArchMemoryMapping},
  breaklines=true, breakatwhitespace=true, postbreak=\mbox{\textcolor{headgray}{$\hookrightarrow$}\space}}
\renewcommand{\file}[1]{\texttt{#1}}   % inherit the size of the surrounding text (tables use \footnotesize)

\fancyhead[L]{\small\textcolor{headgray}{A1-10 shared page -- the short variant}}
\fancyhead[R]{}
\newcommand{\step}[2]{\par\medskip\noindent{\large\bfseries #1}\hfill{\file{#2}}\par\nopagebreak\smallskip}

\begin{document}

{\centering {\LARGE\bfseries A1-10 shared page -- the short variant\par}\vspace{6pt}}

\begin{warnbox}
\textbf{Not tested yet.} This is a shorter version of the tested solution (\file{solution.patch},
explained in \file{walkthrough.pdf}): no new class, no new files, one extra lock. Run \code{task test}
-- done when \textbf{9 of 9} pass. If it does not, compare with \file{minimal.pdf} (tested).
\end{warnbox}

\section*{The idea}

\begin{itemize}
  \item Two processes that attach the same key get the \textbf{same physical page} in their page tables.
  \item Two global arrays remember, per key, \textbf{which physical page} and \textbf{how many processes} use it.
  \item Free the page only when the count reaches \textbf{0}; unmapping a shared page never frees it.
\end{itemize}

\section*{What is shorter than the tested solution}

\begin{tabular}{@{}L{3.6cm} L{5.6cm} L{6.2cm}@{}}
\toprule
 & tested solution & short variant\\\midrule
the table & class \code{SharedPages}, 2 new files & 2 arrays + 1 lock pointer in \file{UserProcess.cpp}\\
locks & two locks, never both at once & both at once, \textbf{always in the same order}\\
lines of logic & $\approx$ 90 & $\approx$ 45\\
\bottomrule
\end{tabular}

\smallskip
\textbf{The one rule that keeps it deadlock-free:} wherever both locks are needed, take
\code{shm\_lock} \textbf{first}, then \code{thread\_lock\_}. Nothing else in the kernel takes \code{shm\_lock}, so
no thread can ever hold \code{thread\_lock\_} and wait for \code{shm\_lock}.

% =====================================================================================
\step{1 -- Plumbing}{syscall-definitions.h · nonstd.h · nonstd.c · Syscall.cpp}

\begin{lstlisting}[style=cpp]
// syscall-definitions.h
#define sc_shm_attach 1680
#define sc_shm_detach 1681

// nonstd.h  (needs #include "types.h" at the top)
extern void* shm_attach(size_t key);
extern int shm_detach(size_t key);

// nonstd.c
void* shm_attach(size_t key)
{
  return (void*) __syscall(sc_shm_attach, key, 0x00, 0x00, 0x00, 0x00);
}
int shm_detach(size_t key)
{
  return __syscall(sc_shm_detach, key, 0x00, 0x00, 0x00, 0x00);
}

// Syscall.cpp, in the switch of syscallException
case sc_shm_attach:
  return_value = currentThread->getProcess()->shmAttach(arg1);
  break;
case sc_shm_detach:
  return_value = currentThread->getProcess()->shmDetach(arg1);
  break;
\end{lstlisting}

% =====================================================================================
\step{2 -- Unmap without freeing}{ArchMemory.h · ArchMemory.cpp}

\begin{lstlisting}[style=cpp]
// ArchMemory.h
bool unmapPage(uint64 virtual_page, bool free_physical_page = true);

// ArchMemory.cpp
bool ArchMemory::unmapPage(uint64 virtual_page, bool free_physical_page)
  ...
  if (free_physical_page)
    PageManager::instance()->freePPN(m.page_ppn);
\end{lstlisting}
\textbf{Why:} the physical page still belongs to the other processes.

% =====================================================================================
\step{3 -- Constants and per-process state}{UserProcess.h}

\begin{lstlisting}[style=cpp]
// at the top, next to the other #defines
#define SHM_KEYS 16                        // keys 1..15 (0 is unused)
#define SHM_BASE 0x0000400000000000ULL     // key k lives at SHM_BASE + k pages

// public:
size_t shmAttach(size_t key);
size_t shmDetach(size_t key);

// private, at the bottom:
bool shm_attached_[SHM_KEYS] = {};          // which keys THIS process has mapped
\end{lstlisting}

% =====================================================================================
\step{4 -- The table: two arrays and a lock}{UserProcess.cpp (top) · main.cpp}

\begin{lstlisting}[style=cpp]
// UserProcess.cpp, below the #includes (add #include "ScopeLock.h" if it is missing)
size_t shm_ppn[SHM_KEYS] = {};     // key -> physical page, 0 = none yet
size_t shm_refs[SHM_KEYS] = {};    // key -> how many processes have it attached
Mutex* shm_lock = nullptr;         // protects both arrays; created in startup()
\end{lstlisting}
\begin{itemize}
  \item Plain arrays of numbers need no constructor -- they are simply zero at boot.
  \item A \code{Mutex} \emph{does} need its constructor, and SWEB never runs constructors of global objects
        -- so only a pointer here, the object is created at boot:
\end{itemize}
\begin{lstlisting}[style=cpp]
// main.cpp: #include "Mutex.h" at the top, and in startup() after Scheduler::instance();
extern Mutex* shm_lock;
shm_lock = new Mutex("shm_lock");
\end{lstlisting}
(\code{extern Mutex* shm\_lock;} goes \emph{above} \code{startup()}, at file level: ``this variable exists in another file''.)

% =====================================================================================
\step{5 -- Attach}{UserProcess.cpp}

\begin{lstlisting}[style=cpp,numbers=left,numberstyle=\tiny\color{headgray}]
size_t UserProcess::shmAttach(size_t key)
{
  if (key == 0 || key >= SHM_KEYS)
    return 0;
  size_t vpn = SHM_BASE / PAGE_SIZE + key;
  ScopeLock global(*shm_lock);          // 1st lock
  ScopeLock local(thread_lock_);        // 2nd lock - always in this order
  if (!shm_attached_[key])
  {
    if (shm_refs[key] == 0)
      shm_ppn[key] = PageManager::instance()->allocPPN();   // first user: new zeroed page
    shm_refs[key]++;
    bool mapped = arch_memory_->mapPage(vpn, shm_ppn[key], 1);
    assert(mapped);
    shm_attached_[key] = true;
  }
  return vpn * PAGE_SIZE;               // attached twice: same address, counted once
}
\end{lstlisting}

% =====================================================================================
\step{6 -- Detach}{UserProcess.cpp}

\begin{lstlisting}[style=cpp,numbers=left,numberstyle=\tiny\color{headgray}]
size_t UserProcess::shmDetach(size_t key)
{
  if (key == 0 || key >= SHM_KEYS)
    return (size_t) -1;
  ScopeLock global(*shm_lock);
  ScopeLock local(thread_lock_);
  if (!shm_attached_[key])
    return (size_t) -1;
  arch_memory_->unmapPage(SHM_BASE / PAGE_SIZE + key, false);   // keep the physical page
  ArchMemory::flushTlb();
  shm_attached_[key] = false;
  shm_refs[key]--;
  if (shm_refs[key] == 0)
    PageManager::instance()->freePPN(shm_ppn[key]);             // last user: free it
  return 0;
}
\end{lstlisting}

% =====================================================================================
\step{7 -- Clean up when a process ends}{UserProcess.cpp, \textasciitilde UserProcess}

Between \code{delete loader\_;} and \code{delete arch\_memory\_;}:
\begin{lstlisting}[style=cpp]
{
  ScopeLock global(*shm_lock);
  for (size_t key = 1; key < SHM_KEYS; ++key)
  {
    if (!shm_attached_[key])
      continue;
    arch_memory_->unmapPage(SHM_BASE / PAGE_SIZE + key, false);
    shm_refs[key]--;
    if (shm_refs[key] == 0)
      PageManager::instance()->freePPN(shm_ppn[key]);
  }
}
\end{lstlisting}
\begin{itemize}
  \item \textbf{Why:} \code{delete arch\_memory\_} frees every mapped page -- also a shared page the other process
        still uses, if this one ends without detaching.
  \item The extra \code{\{ \}} ends the \code{ScopeLock} before \code{delete arch\_memory\_}.
  \item No \code{thread\_lock\_} here: the process is being deleted, none of its threads runs any more.
\end{itemize}

\begin{keybox}[title={Remember}]
\begin{enumerate}
  \item Count the users per key; free the page only at 0.
  \item Unmapping a shared page never frees it (\code{false}).
  \item Two locks: always \code{shm\_lock} first, then \code{thread\_lock\_} -- the same order everywhere.
\end{enumerate}
\end{keybox}

\end{document}
